\expandafter\gdef\csname SNside1\endcsname{\node[anchor=north west,inner sep=0pt,text width=48mm,font=\SideFont,align=left] at ([xshift=138mm,yshift=-583.68pt]current page.north west) {\textcolor{SNink}{\textbf{不变量 · 先数状态再取对数}}\par\vspace{1.5mm}一个码元若有16种可区分状态，它可承载4 bit信息；不是16 bit。码元速率是单位时间的符号数，状态数是每个符号的选择数，把这两个维度分开后，调制与冗余编码就不容易混算。};}
\expandafter\gdef\csname SNside2\endcsname{\node[anchor=north west,inner sep=0pt,text width=48mm,font=\SideFont,align=left] at ([xshift=138mm,yshift=-221.60pt]current page.north west) {\textcolor{SNink}{\textbf{最小反例 · Baud不等于bit/s}}\par\vspace{1.5mm}同为1000 Baud，二进制每码元1 bit得到1 kb/s，16态每码元4 bit得到4 kb/s。两个系统码元率相同、数据率不同，已经足以排除``波特率总等于比特率''的说法。\par\smallskip\begin{tikzpicture}[x=1mm,y=1mm,>=Stealth,node distance=5mm]
\node[draw=SNink!50,fill=SNblue,rounded corners=1mm,text width=34mm,align=center,inner sep=1.2mm,font=\fontsize{7.5}{9}\selectfont] (m0) at (0,0) {1kBaud};
\node[draw=SNink!50,fill=SNblue,rounded corners=1mm,text width=34mm,align=center,inner sep=1.2mm,font=\fontsize{7.5}{9}\selectfont] (m1) at (0,-11) {2态1kb};
\node[draw=SNink!50,fill=SNblue,rounded corners=1mm,text width=34mm,align=center,inner sep=1.2mm,font=\fontsize{7.5}{9}\selectfont] (m2) at (0,-22) {16态4kb};
\draw[->,SNink] (m0) -- (m1);
\draw[->,SNink] (m0) edge[bend left=40] (m2);
\end{tikzpicture}
\par{\fontsize{7.5}{9}\selectfont 同一码元率可以有不同比特率}};}
\expandafter\gdef\csname SNside3\endcsname{\node[anchor=north west,inner sep=0pt,text width=48mm,font=\SideFont,align=left] at ([xshift=138mm,yshift=-95.00pt]current page.north west) {\textcolor{SNink}{\textbf{错项排除 · 分贝不能直接入对数}}\par\vspace{1.5mm}30 dB是功率比1000的对数表示，香农式需要无单位的1000。如果直接把30填入1+S/N，错误发生在输入量的含义，后续对数运算再准确也无法补救。};}
\expandafter\gdef\csname SNside4\endcsname{\node[anchor=north west,inner sep=0pt,text width=48mm,font=\SideFont,align=left] at ([xshift=138mm,yshift=-95.00pt]current page.north west) {\textcolor{SNink}{\textbf{不变量 · 模型决定式子的边界}}\par\vspace{1.5mm}奈氏式限制理想无噪声低通信道的码元率，香农式给出含噪声信道的可靠传输容量。题目同时限定调制阶数时还要保留固定状态数约束，不能把两个上界相加当作新的速率。};}
\expandafter\gdef\csname SNside5\endcsname{\node[anchor=north west,inner sep=0pt,text width=48mm,font=\SideFont,align=left] at ([xshift=138mm,yshift=-95.00pt]current page.north west) {\textcolor{SNink}{\textbf{最小反例 · 中点不能解所有编码}}\par\vspace{1.5mm}差分曼彻斯特中每位中点都会跳变，只看中点方向无法按一般规则恢复比特；起点是否跳变才携带差分编码的信息。第一位还需要前一电平，这正是读波形时应补齐的初始条件。};}
\expandafter\gdef\csname SNside6\endcsname{\node[anchor=north west,inner sep=0pt,text width=48mm,font=\SideFont,align=left] at ([xshift=138mm,yshift=-95.00pt]current page.north west) {\textcolor{SNink}{\textbf{不变量 · 同是引脚也要分维度}}\par\vspace{1.5mm}引脚位于何处属于机械特性；引脚承担请求发送还是接收准备属于功能特性；电压范围属于电气特性。按``这个描述回答什么问题''归类，比按``引脚''这个共同名词归类更稳。};}
\expandafter\gdef\csname SNside7\endcsname{\node[anchor=north west,inner sep=0pt,text width=48mm,font=\SideFont,align=left] at ([xshift=138mm,yshift=-95.00pt]current page.north west) {\textcolor{SNink}{\textbf{最小反例 · 多端口不创造带宽}}\par\vspace{1.5mm}三站连接同一个10 Mb/s集线器，新增端口并未给每站创造独立10 Mb/s通道；一个端口的信号仍占用共享介质。交换机可形成独立链路，但这属于不同设备的行为。};}
\AddToShipoutPictureFG{\begin{tikzpicture}[remember picture,overlay]\ifcsname SNside\thepage\endcsname\csname SNside\thepage\endcsname\fi\end{tikzpicture}}